\subsection{Two elementary moment domains}\label{ms:moment-domain}

The moment inequalities below are used before any finite-dimensional
relaxation. They hold for the actual marginal, without a bounded-support
assumption. In particular, they are unaffected by adding a very small
atom far from the origin.

\begin{theorem}[Signed and absolute third moments]\label{ms:sharp-third}
If $\E X=0$, $\E X^2=1$, $\rho=\E|X|^3<\infty$, and
$\gamma=\E X^3$, then
\[
 \rho\ge\frac{\gamma^2+2}{\sqrt{\gamma^2+4}}.
\]
\end{theorem}
This is the sharp moment relation of Shevtsova
\cite{shevtsova2014complex}; in the cited arXiv version its equivalent
form is Theorem~2. We include a factorization proof to specify the
domain used by the finite certificates.
\begin{proof}
Set $s=\sqrt{\gamma^2+4}$, $a=(s-\gamma)/2$, $b=(s+\gamma)/2$ and
$c=\gamma(\gamma^2+6)/s^3$. Then $a,b>0$, $ab=1$ and $a+b=s$.
Define
\[
 w(y)=|y|^3-cy^3-\frac{12y^2}{s^3}
                          +\frac{6\gamma y}{s^3}+\frac4{s^3}.
\]
For $y,t\ge0$, direct multiplication gives
\[
 w(y)=\frac{2a^2}{s^3}(y-b)^2\{2+(a+3b)y\},\qquad
 w(-t)=\frac{2b^2}{s^3}(t-a)^2\{2+(3a+b)t\}.
\]
Thus $w\ge0$ on the whole real line. Taking expectations and using the
first three signed moments gives
\[
 0\le\E w(X)=\rho-c\gamma-8/s^3
            =\rho-(\gamma^2+2)/s.
\]
All terms are integrable under the stated third-moment hypothesis.
\end{proof}

\begin{lemma}[The signed second moment]\label{ms:sign-moment}
With the same normalization and $\beta=\E[X|X|]$,
\[
 |\beta|<1,\qquad
 \rho\ge\frac{1+\beta^2}{\sqrt{1-\beta^2}}.
\]
\end{lemma}
\begin{proof}
Both signs have positive probability. Put
$m=\E X_+=\E X_->0$ and $a_\pm=\E X_\pm^2=(1\pm\beta)/2>0$.
If $p_\pm=\Pp(\pm X>0)$, Cauchy's inequality gives
$m^2\le p_\pm a_\pm$. Since $p_++p_-\le1$, this implies
\[
 m^2(1/a_++1/a_-)\le1,
 \qquad 2m\le\sqrt{1-\beta^2}.
\]
Another application of Cauchy's inequality, to
$X_\pm^{1/2}$ and $X_\pm^{3/2}$, gives
$m\E X_\pm^3\ge a_\pm^2$. Summing yields
$m\rho\ge(1+\beta^2)/2$, which proves the claim. A possible zero atom
enters only through $p_++p_-\le1$ and causes no exception.
\end{proof}
